\documentclass[10pt,a4paper]{amsart}
\usepackage[no-math]{fontspec}
\usepackage{amsmath,amssymb,mathtools,xfrac,xcolor}
\usepackage[a4paper,margin=22mm]{geometry}
\setmainfont{Latin Modern Roman}[SmallCapsFont=lmromancaps10-regular.otf]
\usepackage[colorlinks=true,linkcolor=blue,citecolor=teal,urlcolor=blue]{hyperref}
\sloppy\newcommand{\ie}{i.e.\ }\newcommand{\resp}{resp.\ }\newcommand{\skpt}{\par\smallskip}
\usepackage[scr=rsfs,cal=euler]{mathalfa}
\multlinegap0pt
\hyphenation{exam-ple sequen-ce}
\renewcommand{\theenumi}{\roman{enumi}}
\newcommand{\psfrac}[2]{\sfrac{(#1)}{#2}}
\newcommand{\sbullet}{{\scriptscriptstyle\bullet}}
\newcommand{\csbullet}{{\raisebox{1pt}{$\sbullet$}}}
\newcommand{\bA}{\boldsymbol{A}}
\newcommand{\bB}{\boldsymbol{B}}
\newcommand{\bE}{\boldsymbol{E}}
\newcommand{\bme}{\boldsymbol{e}}
\newcommand{\bmf}{\boldsymbol{f}}

\usepackage{bbm}
\usepackage{mathtools}

\theoremstyle{plain}
\newtheorem{theorem}{Theorem}[section]
\newtheorem{lemma}[theorem]{Lemma}
\newtheorem{proposition}[theorem]{Proposition}

\theoremstyle{definition}
\newtheorem{definition}[theorem]{Definition}
\newtheorem{remark}[theorem]{Remark}

\let\con\overline
\let\mc\mathcal
\let\op\operatorname

\newcommand{\N}{\mathbb{N}}
\newcommand{\R}{\mathbb{R}}
\newcommand{\Z}{\mathbb{Z}}
\newcommand{\HH}{\mathbb{H}}

\let\al\alpha
\let\ga\gamma
\let\eps\varepsilon
\let\la\lambda
\let\Si\Sigma
\let\Om\Omega

\newcommand{\dd}{\mathrm{d}}

\DeclareMathOperator{\comp}{c}
\DeclareMathOperator{\Ell}{ell}
\DeclareMathOperator{\Op}{Op}
\DeclareMathOperator{\ran}{ran}
\DeclareMathOperator{\supp}{supp}
\DeclareMathOperator{\vol}{vol}

\newcommand{\bigo}{\mathcal{O}}

\begin{document}
\title[Low-energy magnetic eigenstates]{Realization of invariant measures by low-energy magnetic Laplace eigenstates}
\author{Laurent Charles\and Thibault Lefeuvre}
\date{}\maketitle\setcounter{section}{1}
Let $(\Sigma,g)$ be a closed connected oriented hyperbolic (constant curvature $-1$) surface of genus $\textsl{g} \geq 2$. Let $L \to \Sigma$ be a Hermitian line bundle equipped with a unitary connection $\nabla$, and denote by $F_\nabla = -i B \vol \in C^\infty(\Sigma, \Lambda^2 T^*\Sigma)$ the curvature $2$-form of $\nabla$, where $B \in C^\infty(\Sigma)$, and $\vol$ is the Riemannian volume. We~call~$B$ the \emph{magnetic field}. The \emph{magnetic Laplacian} is defined as
\begin{equation}
\label{equation:delta}
\Delta_L := \tfrac{1}{2} \nabla^* \nabla : C^\infty(\Sigma,L) \to C^\infty(\Sigma,L).
\end{equation}
More generally, taking $L^{\otimes k}$ for $k \in \Z_{\geq 0}$ and the induced connection $\nabla^{\otimes k}$, one can form similarly to \eqref{equation:delta} a Laplacian $\Delta_k$ acting on $C^\infty(\Sigma,L^{\otimes k})$. {\color{black}Note that the curvature on $L^{\otimes k}$ is $-ikB\vol$.} The operator $k^{-2}\Delta_k$ is a twisted semiclassical (pseudo)differential operator, with semiclassical parameter $h := k^{-1} > 0$.

Throughout this article, we~will further assume that $B$ is \emph{constant}. Note that, by~Gauss-Bonnet, this implies that $2B(\textsl{g}-1) \in \Z$. The purpose of this paper is to study the semiclassical limits of the Laplace eigenstates
\begin{equation}
\label{equation:eigenstates}
k^{-2}\Delta_k u_k = (E+\eps_k)u_k, \qquad u_k \in C^\infty(\Sigma,L^{\otimes
k}), \; \|u_k\|_{L^2(\Sigma,L^{\otimes k})} = 1,
\end{equation}
in the regime $k \to +\infty$, where $E \geq 0$ and $\eps_k \to_{k \to +\infty} 0$. {\color{black}This means that we consider the regime where the magnetic strength $kB$ (that is, the curvature of $L^{\otimes k}$) tends to infinity.} We say that $u_k$ converges to the semiclassical defect measure $\mu$ (defined on~$T^*\Sigma$) if for all $a \in C^\infty(T^*\Sigma)$,
\begin{equation}
\label{equation:convergence}
\langle\Op_{k}(a)u_{k},u_{k}\rangle_{L^2} \underset{k \to \infty}\to \int_{T^*\Sigma} a(x,\xi) \dd\mu(x,\xi),
\end{equation}
where $\Op$ is a certain semiclassical magnetic quantization on $\Sigma$ (see Section~\ref{section:microlocal} for a definition). We~denote by $u_k \rightharpoonup_{k \to \infty} \mu$ this convergence. By the normalization assumption $\|u_k\|_{L^2(\Sigma,L^{\otimes k})} = 1$, $\mu$ is a probability measure; in addition, $\mu$ must be supported on the compact energy shell $\{p=E\} \subset T^*\Sigma$ (where $p(x,\xi) := \tfrac{1}{2}|\xi|^2$). The limit \eqref{equation:convergence} may not exist as $k \to +\infty$; however, it~always exists along a subsequence $(k_n)_{n \geq 0}$ and we will often drop the subsequence notation. We~shall see that three different semiclassical regimes appear according to the value of $E$, affecting the possible behaviour of $\mu$.

\subsection{Main results}

Let $\omega_0$ be the Liouville symplectic $2$-form on $T^*\Sigma$. Set
\[
\Omega := \omega_0 + i \pi^* F_\nabla,
\]
where $\pi : T^*\Sigma \to \Sigma$ is the projection. This is the Liouville $2$-form with a magnetic correction; it is still a symplectic $2$-form on $T^*\Sigma$. The (semiclassical) principal symbol of $k^{-2}\Delta_k$ is $p(x,\xi) = \tfrac{1}{2}|\xi|^2_g$. The Hamiltonian vector field $H_p^\Omega$ of $p$ computed with respect to $\Omega$ is defined through the relation
\[
dp(\csbullet) = \Omega(\csbullet, H^{\Omega}_p).
\]
The Hamiltonian flow $(\Phi_t)_{t \in \R}$ generated by $H_p^{\Omega}$ is called the \emph{magnetic flow}. As any autonomous Hamiltonian flow, $(\Phi_t)_{t \in \R}$ preserves the (compact) energy layers $\{p=E\}$ and a natural smooth probability measure $\mu_{\mathrm{Liouv}}$ on $\{p=E\}$ called the Liouville measure (see Section~\ref{ssection:geo} for these definitions).

Let $S\Sigma \to \Sigma$ be the unit tangent bundle. Denote by $(\varphi_t)_{t \in \R}$ the (Anosov) geodesic flow, $(R_t)_{t \in \R}$ the $2\pi$-periodic rotation in the circle fibers of $S\Sigma$, and $(h_t)_{t \in \R}$ the stable horocyclic flow. Define
\[
E_c := \tfrac{1}{2}B^2, \quad T_E = (B^2-2E)^{-1/2} \text{ if } E < E_c, \quad T_E := (2E-B^2)^{-1/2} \text{ if } E > E_c.
\]
The energy $E_c$ is called the \emph{critical energy}. The following fact is
well-known \cite{Arnold-61,Ginzburg-96} and will be reproved quickly in
Section~\ref{ssection:magnetic-flow}. For $E > 0$, the magnetic flow $(\Phi_t)_{t \in \R}$ on $\{p=E\}
\subset T^*\Sigma$ is conjugate either to:
\begin{enumerate}

\item the reparametrized rotation flow $(R_{t/T_E})_{t \in \R}$ of period
$2\pi \cdot T_E$ if $E < E_c$ (elliptic case),

\item the horocyclic flow $(h_t)_{t \in \R}$ if $E=E_c$ (parabolic case),

\item the reparametrized geodesic flow $(\varphi_{t/T_E})_{t \in \R}$ if $E>E_c$ (hyperbolic case).
\end{enumerate}
Notice that for $E < E_c$ the space of orbits at energy $E$ is diffeomorphic to $\Si$ itself, so~the $(\Phi_t)_{t \in \R}$ invariant probability measures of $\{ p = E\}$
identify with the probability measures of $\Si$. This holds as well for
$E=0$ because the flow is stationary on $\{ p =0 \}$.

We will prove that this transition in the dynamics of the magnetic flow at the critical energy $E=E_c$ translates at the quantum level into the following result on semiclassical defect measures:

\begin{theorem}[The three classical/quantum regimes]
\label{theorem:main}
The following holds under the above assumptions on $(\Sigma,g)$ and $(L,\nabla)$:
\begin{enumerate}

\item\label{theorem:maini} \emph{Low energies regime.} If $0 \leq E < E_c$, then for any $(\Phi_t)_{t \in \R}$ invariant probability measure $\mu$ on $\{p=E\}$, there exists a sequence $(u_{k})_{k \geq 0}$ satisfying \eqref{equation:eigenstates} such that $u_{k} \rightharpoonup_{k \to \infty} \mu$;

\item\label{theorem:mainii} \emph{Critical energy regime.} For any
sequence $(u_k)_{k \geq 0}$ satisfying \eqref{equation:eigenstates} with $E=E_c$, $u_{k} \rightharpoonup_{k \to \infty} \mu_{\mathrm{Liouv}}$

\item\label{theorem:mainiii} \emph{High energies regime.} If $E_c < a < b$, consider for
all $k \geq 0$, the set of eigenstates $(u_{k,j})_{k,j \in A}$ such that
$k^{-2} \Delta_k u_{k,j} = \lambda_{k,j} u_{k,j}$ with $\lambda_{k,j} \in
[a,b]$ and $A \subset \Z_{\geq 0}^2$. Then there exists a density one subset
$A_\star \subset A$ such that for all sequences $(u_{k_n,j_n})_{n \geq 0}$
with $(k_n,j_n) \in A_\star$ and $k_n \to \infty$, $u_{k_n,j_n} \rightharpoonup_{n \to \infty} \mu_{\mathrm{Liouv}}$.
\end{enumerate}
\end{theorem}
\section{Background}

\label{section:dynamical-geometric-background}

Throughout this section, we~let $(\Sigma, g)$ denote a closed, connected, oriented hyperbolic surface—that is, a Riemannian surface with constant curvature $-1$. We~establish the following results:

\begin{itemize}
\item In Section~\ref{ssection:geo}, we~recall several technical results concerning the geometry and dynamics on the tangent bundle $T\Sigma$. In~particular, we~emphasize the unique ergodicity of the horocyclic flow, a key ingredient in the proof of Theorem~\href{https://jep.centre-mersenne.org/item/10.5802/jep.333.pdf#page=5}{1.2}.

\item In Section~\ref{ssection:magnetic-flow}, we~introduce the magnetic flows on both the tangent and cotangent bundles, and show that they are equivalent under the identification induced by the metric. We~also prove a technical result on the propagation of functions along the magnetic flow, which will be instrumental in the proof of Theorem~\href{https://jep.centre-mersenne.org/item/10.5802/jep.333.pdf#page=5}{1.2}.

\item Finally, in Section~\ref{ssection:landau}, we~review standard properties of magnetic Laplacians. In~particular, we~explicitly compute the bottom of their spectrum—the so-called Landau levels. \end{itemize}

\subsection{Geometry and dynamics on the tangent bundle}

\label{ssection:geo}
Let $T\Sigma$ be the tangent bundle of $\Sigma$ and
\begin{equation}
\label{equation:footpoint}
\pi : T\Sigma \to \Sigma
\end{equation}
be the footpoint projection. We~refer the reader to \cite[Ch.\,1]{Paternain-99} or \cite[Ch.\,13]{Lefeuvre-book} for background material on the geodesic flow.

\subsubsection{Structural equations} The geodesic flow $\varphi_t : T\Sigma \to T\Sigma$ is defined as
\[
\varphi_t(x,v) := (\gamma(t), \dot{\gamma}(t)),
\]
where $\gamma : \R \to \Sigma$ is the (unique) curve solving Newton's equation:
\begin{gather} \label{eq:newton_geodesic}
\nabla_{\dot{\gamma}(t)}\dot{\gamma}(t) = 0, \qquad \gamma(0)=x, \; \dot{\gamma}(0)=v,
\end{gather}
and $\nabla$ is the Levi-Civita connection. We~let
\[
X := \partial_t\varphi_t|_{t=0} \in C^\infty(T\Sigma, T(T\Sigma))
\]
be its generator.

As $\Sigma$ is oriented, there is a well-defined fiberwise rotation $R_\theta : T\Sigma \to T\Sigma$ by angle $\theta \in [0,2\pi)$ in the fibers of $T\Sigma$. Let $V$ be the generator of this rotation, namely
\[
V := \partial_\theta R_\theta|_{\theta = 0}.
\]
Let $V_\perp$ be the Euler vector field of $T\Sigma$, i.e. the generator of the flow
\begin{equation}
\label{equation:flot-vperp}
\varphi_t^{V_\perp}(x,v) = (x,e^tv).
\end{equation}
Finally, define
\[
X_\perp := [V,X].
\]
The vector fields $\{X,X_\perp,V,V_\perp\}$ form a basis of $T(T\Sigma)$ at any point of $T\Sigma$. The metric for which this is an orthonormal frame is called the \emph{Sasaki metric}.

We will need the following result:

\begin{lemma}
The above vector fields satisfy the commutation relations:
\begin{equation}
\label{equation:commutator-relations}
\begin{aligned}
[X,X_\perp] &= - |v|^2 V,& [X,V] &=-X_\perp, & [X_\perp,V] &= X, \\
[V_\perp,X] &= X, & [V_\perp, X_\perp] &= X_\perp, & [V,V_\perp] &= 0.
\end{aligned}
\end{equation}
\end{lemma}

We refer to \cite[Lem.\,15.2.1]{Lefeuvre-book} for a proof. (In this reference, $X_\perp$ is denoted by~$H$; the computations are done on the unit tangent bundle but the above relations follow easily by a scaling argument.) Notice that in our case, the sectional curvature is $\kappa = -1$.

\subsubsection{Horocyclic flow} For $\lambda \geq 0$, define
\[
S^\lambda\Sigma := \{|v|_g=\lambda\}.
\]
We let $S\Sigma := S^{1}\Sigma$ be the unit tangent bundle. It is straightforward to verify that $X,X_\perp$ and $V$ preserve the layers $S^\lambda\Sigma$.

The geodesic flow $(\varphi_t)_{t \in \R}$ is Anosov on $S\Sigma$ in the sense that {\color{black}there exists a continuous flow-invariant decomposition}
\[
T(S \Sigma) = \R X \oplus E_s \oplus E_u,
\]
where
\[
\begin{split}
|d\varphi_t(w)| \leq e^{-t} |w|, &\qquad \forall t \geq 0,\; w \in E_s \\
|d\varphi_{-t}(w)| \leq e^{-t} |w|, &\qquad \forall t \geq 0,\; w \in E_u.
\end{split}
\]
Here $|\csbullet|$ denotes the norm induced by the \emph{Sasaki metric} on $T\Sigma$, which is defined as the metric for which $\{X,X_\perp,V,V_\perp\}$ is an orthonormal frame.

We define the (stable) horocyclic vector field by
\[
U_+ := X_\perp - V.
\]
Let $(h_t)_{t \in \R}$ be the (stable) \emph{horocyclic flow} generated by $U_+$ on $S\Sigma$. Recall that the Liouville measure on $S\Sigma$ is the unique volume form $\mu_{\mathrm{Liouv}}$ such that $\mu_{\mathrm{Liouv}}(X,X_\perp,V) = \mathrm{cst}$. The constant is normalized such that $\mu_{\mathrm{Liouv}}$ is a probability measure on $S\Sigma$. The Liouville measure is invariant by $X, X_\perp$ and $V$; it is thus also invariant by $U_+$.

\begin{theorem}
\label{lemma:horocycle}
The following holds:
\begin{enumerate}
\item\label{lemma:horocyclei} The flow $(h_t)_{t \in \R}$ is uniquely ergodic on $S\Sigma$ and $\mu_{\mathrm{Liouv}}$ is the only flow-invariant probability measure.

\item\label{lemma:horocycleii} Let $\theta\in(0,1/2)$ be such that $\theta(1-\theta) \leq \lambda_1(\Sigma)$, where $\lambda_1(\Sigma) > 0$ is the first non-zero eigenvalue of the Laplacian $\Delta_g$ on functions. Then there exists $C > 0$ such that for all $a \in C^{\infty}(S\Sigma)$, for all $T > 0$,
\[
\sup_{v \in S\Sigma}\, \biggl|\dfrac{1}{T}\int_0^T a(h_t(v))~ \dd t - \int_{S\Sigma} a(v) ~\dd\mu_{\mathrm{Liouv}}\biggr| \leq C T^{-\theta} \|a\|_{H^3(S\Sigma)},
\]
where $H^3(S\Sigma)$ denotes the $L^2$-based Sobolev norm of order $3$.
\end{enumerate}
\end{theorem}

Unique ergodicity in constant curvature was established by Furstenberg \cite{Furstenberg-73}, while the rate of convergence of ergodic averages was proved by Burger \cite{Burger-90}. A~refinement of this result can be found in \cite{Flaminio-Forni-03}.

\subsection{Magnetic hyperbolic flow}

\label{ssection:magnetic-flow}

In this paragraph, we~discuss the magnetic flow both on the tangent and the cotangent bundles.

\subsubsection{Magnetic flow on the tangent bundle}
\label{sssection:magnetic-tsigma}
A {\em magnetic geodesic}
is a curve $\ga $ of $\Sigma$ which satisfies the equation
\begin{equation} \label{eq:Newton_Lorentz}
\nabla_{\dot{\gamma}(t)} \dot{\gamma}(t) = - B j_{\gamma(t)}\dot{\ga} (t),
\qquad \gamma(0)=x,\; \dot{\gamma}(0)=v,
\end{equation}
where $(x,v) \in T\Sigma$, $\nabla$ stands for the Levi-Civita covariant derivative, $j$ is the almost-complex structure and
$ B \in C^\infty(\Sigma)$ is the magnetic {\em intensity}. {\color{black}As $\Sigma$ is oriented, and equipped with a metric $g$, the almost-complex structure $j \in C^\infty(\Sigma,\mathrm{End}(T\Sigma))$ is the fiberwise endomorphism of $T\Sigma$ given by rotation of tangent vectors by an angle $+\pi/2$.} In the sequel we assume that $B$ is constant
and positive.
The {\em magnetic flow} $(\Phi_t)_{t \in \R}$ is the autonomous flow of $T \Si$ such that for
any curve $\ga$ satisfying \eqref{eq:Newton_Lorentz},
\[
\Phi_t (x,v) := (\ga (t), \dot {\ga} (t)), \qquad \forall t \in \R.
\]
\begin{lemma} The magnetic flow is generated by the vector field $F := X - B V$.
\end{lemma}

\begin{proof} We work in a chart $U$ of $\Si$ that we identify with an open
subset of $\R^2$, so~$TU = U \times \R^2$. Let $(x,v) \in TU$. Let $c_1$,
$c_2$ be a geodesic and a magnetic geodesic respectively such that $ c_1
(0) = x = c_2 (0)$, $\dot{c}_1 (0) =v = \dot{c}_2 (0)$. Then $X(x,v) = (
\dot{c}_1 (0), \ddot{c}_1 (0))$ and $F(x,v) = (
\dot{c}_2 (0), \ddot{c}_2 (0)) $, thus
\[
F(x,v) - X(x,v) = (0, \ddot{c}_2 (0) - \ddot{c}_1 (0)) = (0,
(\nabla_{\dot{c}_2} \dot{c}_2) (0) - (\nabla_{\dot{c}_1} \dot{c}_1) (0)
) = (0, -B j_{x} v),
\]
where we have used that $\nabla_{\dot{c}_i} \dot{c}_i = \ddot{c}_i +
\Gamma_{c_i} (\dot{c}_i) (\dot{c}_i)$, $\Gamma$ being the connection
one-form, and then Newton's equations \eqref{eq:newton_geodesic},
\eqref{eq:Newton_Lorentz}. To conclude use that $V (x,v) = (0, j_x v)$.
\end{proof}

\begin{remark}The flow $(\Phi_t)_{t \in \R}$ is considered here on the tangent bundle $T\Sigma$, whereas it was initially defined on the cotangent bundle $T^*\Sigma$ in the introduction. In~Section~\ref{sssection:flow-cotangent}, we~show that these two flows are equivalent under the musical isomorphism induced by the metric. To avoid unnecessary notation, unless specified explicitly, we~do not distinguish between the flows on $T\Sigma$ and $T^*\Sigma$ in what follows.
\end{remark}

Recall that $(\varphi_t)_{t \in \R}$ is the geodesic flow (generated by $X$), $(R_t)_{t \in \R}$ is the $2\pi$-periodic rotation in the circle fibers of $S\Sigma$ (generated by $V$).

\begin{proposition}
\label{proposition:dynamic}
For any $\lambda> 0$, the following holds:
\begin{enumerate}

\item\label{proposition:dynamici} If $ \lambda< B$, $(\Phi_t)_{t \in \R}$ on $S^\lambda\Sigma$ is conjugate to $(R_{t/T_\lambda})_{t \in \R}$ on $S\Sigma$, where $T_\lambda = (B^2 - \lambda^2)^{-\sfrac{1}{2}}$;

\item\label{proposition:dynamicii} If $ \lambda = B$, then $(\Phi_t)_{t \in \R}$ on $S^\lambda\Sigma$ is conjugate to $(h_t)_{t \in \R}$ on $S\Sigma$;

\item\label{proposition:dynamiciii} If $ \lambda>B$, then $(\Phi_t)_{t \in \R}$ on $S^\lambda\Sigma$ is conjugate to $(\varphi_{t/T'_\lambda})_{t \in \R}$ on $S\Sigma$, where $ T_\lambda' = (\lambda^2 -B^2)^{-\sfrac{1}{2}}$.
\end{enumerate}
\end{proposition}

\begin{proof}
For $\lambda > 0$, let $\Psi_\lambda : T\Sigma \to T\Sigma$ be the map defined by $\Psi_\lambda(x,v) := (x,\lambda v)$. Observe that $\Psi_\lambda^* X = \lambda X$ and $\Psi_\lambda^*V=V$. The flow of $F = X-BV$ on $S^\lambda \Sigma$ is thus conjugate to the flow of $\Psi_\lambda^* F = \lambda X -BV$ on $S\Sigma$.

{\color{black}As $(\Sigma,g)$ is a hyperbolic surface, there exists a discrete torsion-free subgroup $\Gamma < \mathrm{PSL}(2,\R)$ such that $S\Sigma = \Gamma \backslash S\HH^2 = \Gamma\backslash\mathrm{PSL}(2,\R)$.} The vector field $\lambda X - BV$ is represented by the following matrix element of $\mathfrak{sl}(2,\R)$:
\[
\arraycolsep3.5pt
\lambda \begin{pmatrix} 1/2 & 0 \\ 0 & -1/2 \end{pmatrix} - B \begin{pmatrix} 0 & 1/2 \\ -1/2 & 0\end{pmatrix} = \begin{pmatrix} \lambda/2 & -B/2 \\ B/2 & -\lambda/2 \end{pmatrix}.
\]
For $\lambda = 1, B=0$, this is the generator of the geodesic flow; for $\lambda=0$, $B=1$, this is the generator of the rotation flow in the circle fibers, see \cite[\S 2.1]{Flaminio-Forni-03} for details. The eigenvalues of this matrix are $\pm\tfrac{1}{2}(\lambda^2-B^2)^{1/2}$ if $\lambda \geq B$ and $\pm \tfrac{i}{2}(B^2-\lambda^2)^{1/2}$ if $B \geq \lambda$. Notice that for $\lambda = B$, one finds the matrix
\[
\arraycolsep3.5pt
\lambda/2 \begin{pmatrix} 1 & -1 \\ 1 & -1 \end{pmatrix}
\]
which is conjugate by an element of $\mathrm{PSL}(2,\R)$ to
\[
\arraycolsep3.5pt
\begin{pmatrix} 0 & 1 \\ 0 & 0 \end{pmatrix},
\]
the generator of the horocyclic flow.

If two matrices in $\mathfrak{sl}(2,\R)$ are conjugate (by an element of $\mathrm{PSL}(2,\R)$), then their corresponding flows are conjugate too. Hence, if~$B > \lambda$, $(\Phi_t)_{t \in \R}$ on $S^\lambda\Sigma$ is conjugate to $(R_{t(B^2-\lambda^2)^{1/2}})_{t \in \R} = (R_{t/T_\lambda})_{t \in \R}$ on $S\Sigma$. If $B=\lambda$, $(\Phi_t)_{t \in \R}$ on $S^\lambda\Sigma$ is conjugate to $(h_t)_{t \in \R}$ on $S\Sigma$. If $B < \lambda$, $(\Phi_t)_{t \in \R}$ on $S^\lambda\Sigma$ is conjugate to $(\varphi_{t(\lambda^2-B^2)^{1/2}})_{t \in \R} = (\varphi_{t/T'_\lambda})_{t \in \R}$ on $S\Sigma$.
\end{proof}

\subsubsection{Magnetic flow on the cotangent bundle} \label{sssection:flow-cotangent}

We now define the magnetic flow on the cotangent bundle, and show that it coincides with that on the tangent bundle up to the identification provided by the metric. Let
\[
\flat : T\Sigma \to T^*\Sigma, \qquad \flat(x,v) := (x, g_x(v,\csbullet))
\]
be the musical isomorphism. Let $A$ be the Liouville
form on $T^*\Sigma$ defined by
\begin{equation}
\label{equation:liouville}
A_{(x, \xi) } := \xi (\dd_{x, \xi} \pi (\csbullet)),
\end{equation}
where $\pi : T^*\Sigma \to \Sigma$ denotes the projection. Consider the symplectic form
\[
\Om \in C^\infty(T^* \Si,\Lambda^2 T^*(T^*\Sigma)), \qquad \Om := dA + B~ \pi^*\op{vol},
\]
where $\op{\vol} \in C^\infty(\Si,\Lambda^2 T^*\Sigma)$ the Riemannian volume. Let $p(x,\xi) := \tfrac{1}{2} | \xi|^2$, and define the flow $(\Phi_t^{T^*\Sigma})_{t \in \R}$ as the Hamiltonian flow of the function $p \in C^\infty(T^*\Sigma)$ with respect to the symplectic form $\Omega$. Namely $(\Phi_t^{T^*\Sigma})_{t \in \R}$ is generated by $H^{\Omega_p}$ such that
\[
dp = \Omega(\csbullet, H^\Omega_p).
\]
To avoid confusion, let us denote temporarily $(\Phi_t^{T\Sigma})_{t \in \R}$ the magnetic flow defined on $T\Sigma$ in Section~\ref{sssection:magnetic-tsigma}.

\begin{lemma}
The map $\flat$ intertwines $(\Phi_t^{T\Sigma})_{t \in \R}$ and $(\Phi_t^{T^*\Sigma})_{t \in \R}$, that is
\[
\Phi_t^{T^*\Sigma} \circ \flat = \flat \circ \Phi_t^{T\Sigma}, \qquad \forall t \in \R.
\]
\end{lemma}

\begin{proof}
Define the $1$-form $\alpha := \flat^*A$ on $T\Sigma$. It satisfies
\[
\alpha_{(x,v)} = g_x(\dd_{x,v}\pi(\csbullet),v),
\]
see \cite[Lem.\,1.37, item\,1]{Paternain-99}. Hence
\[
\flat^*\Omega = d\alpha + B~ \pi^*\vol,
\]
where $\pi : T\Sigma \to \Sigma$ is the footpoint projection. Define $h(x,v) := |v|^2_g/2$. Let $H_h^{\flat^*\Omega}$ be the Hamiltonian vector field of $h$, computed with respect to $\flat^*\Omega$ on $T\Sigma$, namely
\[
dh = \flat^*\Omega(\csbullet, H_h^{\flat^*\Omega}).
\]
It is immediate that $\flat^* H_p^{\Omega} = H_h^{\flat^*\Omega}$ so the claim boils down to showing that
\[
H_h^{\flat^*\Omega} = X - B V.
\]
For that, we~compute:
\begin{equation}
\label{equation:mouch}
\flat^*\Omega(\csbullet,X-BV) = \dd\alpha(\csbullet, X) - B \dd \alpha(\csbullet, V) + B \pi^*\vol(\csbullet, X).
\end{equation}
A quick computation reveals that $\dd \alpha(\csbullet, X) = dh$. Indeed, $\dd \alpha(Y,X) = 0 = dh(Y)$ for $Y = X,X_\perp,V$ and $\dd \alpha(V_\perp,X) = |v|^2 = dh(V_\perp)$ using the $2$-homogeneity of $h$ and that $V_\perp$ is the Euler vector field.

In addition, using \cite[Prop.\,1.24]{Paternain-99} for the formula for $\dd \alpha$, we~find:
\[
\dd \alpha(\csbullet,V) = - g_x(j(x)v, \dd \pi(\csbullet)) = \pi^*\vol(\csbullet,X),
\]
so the last two terms in \eqref{equation:mouch} cancel out. This proves that $\flat^*\Omega(\csbullet,X-BV) = dh$, and thus $H_h^{\flat^*\Omega} = X - B V$.
\end{proof}
\subsection{Magnetic hyperbolic Laplacian} \label{ssection:landau} By the Gauss-Bonnet formula, the metric being
hyperbolic, the volume is
\begin{equation}
\int_{\Si} \vol = 4 \pi (\textsl{g} -1),
\end{equation}
where $\textsl{g}\geqslant
2 $ is the genus of $\Si$. Let $L \to \Sigma$ be a Hermitian line bundle with
a Hermitian connection
$\nabla$ whose curvature is $-i B \op{vol}
$, $B > 0$ being a positive constant. Since the Chern class of $L$ is $ (2
\pi)^{-1} B \cdot [ \op{vol}]$, the
degree of $L$ is
\begin{equation}
\label{equation:deg}
\mathrm{deg}(L) = (2 \pi)^{-1} B \int_{\Si} \vol = 2 B (\textsl{g}-1) \in \Z.
\end{equation}
Thus, the pair $(L,
\nabla)$ exists if and only if $2 B (\textsl{g}-1)$ is an integer. When it exists, it
is unique up to tensor product by a flat Hermitian line bundle. The
space of flat Hermitian line bundles up to equivalence is in one-to-one
correspondence trough the holonomy representation with the character space
$\op{Mor} (\pi_1 (\Si), \op{U}(1)) \simeq \op{U} (1)^{2\textsl{g}}$. An~example of a pair $(L, \nabla)$ is provided by the canonical bundle $K := T^{1,0} \Si$,
endowed with its Chern connection, whose curvature is $- i \vol$, so~in this
case $B=1$.

Let $\Delta := \frac{1}{2} \nabla^* \nabla $ be the Laplacian acting on $C^\infty
(\Si, L)$. Since $\Si$ is compact, $\Delta$~has a discrete spectrum and each
eigenvalue has finite multiplicity. We~denote by $\la_0 \leqslant \la_1
\leqslant \cdots $ the eigenvalues of $\Delta$. The following result was
proved in \cite{Iengo-Li-94} (see also \cite{Tejero-06, Landau2}):

\begin{proposition} \label{theo:Landau_Level}
The following holds:
\begin{enumerate}

\item For any pair $(L, \nabla)$ with curvature $F_\nabla=-iB
\op{vol}$ with $B>0$, the first $N := \lfloor B \rfloor $ eigenvalues are
given by
\begin{gather} \label{eq:valeur_propre}
\la_{m } = B (\sfrac{1}{2} + m) - \dfrac{ m (m+1) } {2},\qquad 0
\leqslant m < N,
\end{gather}
and the multiplicity of $\la_m$ is $2 (\textsl{g}-1)(B-\sfrac{1}{2} - m)$ when $m \leqslant N- 2$.

\item If $(L,\nabla) = (K^r,\nabla^{\mathrm{Chern}})$ with $ r \geq 1$, then \eqref{eq:valeur_propre} holds for
$m \leqslant r$ and the remaining part of the spectrum is given by
$\la_{r+n} = \la_r + \frac{1}{2} \mu_n$ where $\{ \mu_n, \; n \geq 0 \}$ is the
spectrum of the Laplace-Beltrami operator of $(\Sigma, g)$.
\end{enumerate}
\end{proposition}

We call the eigenvalues $\la_m$ given by \eqref{eq:valeur_propre} and their eigenspaces the Landau levels.
These quantum levels correspond to the classical energy
levels $S^\lambda \Si$ with $\lambda <B$ on which the magnetic flow is periodic. A~first
clue of this fact is that \eqref{eq:valeur_propre} applied with $m = B$ gives $\frac{1}{2} B^2$ which is exactly the classical
energy $\frac{1}{2} \lambda^2$ at the critical value $\lambda =B$.
Moreover, as was noticed by Comtet \cite{Comtet-86}, \eqref{eq:valeur_propre} can be interpreted as a
Bohr-Sommerfeld condition.

On a different perspective, a common characteristic
of the magnetic flow below the critical energy and the Landau levels is
that they do not depend on the choice of hyperbolic metric $g$. The part of the spectrum above the
critical energy has a different nature (as seen with $L = K^r$ for instance).

\begin{proof}[Sketch of proof] This follows from the Riemann-Roch theorem and the
following two identities
\begin{gather} \label{eq:bochner_bosonic}
\Delta = \Box_L + \tfrac{1}{2} B, \qquad \con{\partial}_L
\con{\partial}_L^* = \Box_{L \otimes K^{-1}} +B -1.
\end{gather}
Here, $L$ is equipped with its holomorphic structure whose Chern connection
is $\nabla$, $\con{\partial}_L: C^{\infty} (L)
\rightarrow C^{\infty} (L \otimes K^{-1})$ is the $d$-bar
operator with the identification $\con{ K} = K^{-1}$ induced by
the metric, and $\Box_L = \con{\partial}_L^*
\con{\partial}_L$. The first identity of \eqref{eq:bochner_bosonic} is a
classical Bochner identity, the second one is proved in \cite[Th.\,7.1]{Landau1}.

For any $m \in \N$, let $\con{\partial}_m = \con{\partial}_{L \otimes
K^{-m}}$ and $\Box_m = \Box_{L \otimes K^{-m}}$. The curvature of $L
\otimes K^{-m}$ being $\frac{1}{i} (B -m) \op{vol}$, the second identity
gives
\[
\con{\partial}_m \Box_m = (
\Box_{m+1} + (B - m) - 1) \con{\partial}_m.\]
Introduce the operator $\Box_{m}^{-1}$
of $C^\infty(L \otimes K^{-m})$ equal to $0$ on $\ker \Box_m$
and inverting $\Box_m$ on the orthogonal of $\ker \Box_m$. Then
$\Box_m^{-1} \con{\partial}_m^*$ is the inverse of $\con{\partial}_m$ on
$(\ker \Box_m)^{\perp}$, so~by the previous equation
\begin{gather*}
\Box_m = \Box_m^{-1} \con{\partial}_m^* \bigl(\Box_{m+1} + B - (m+1)
\bigr) \con{\partial}_m.
\end{gather*}
As a consequence, the spectra satisfy
\begin{gather} \label{eq:rec}
\op{Sp} (\Box_m) \setminus \{ 0 \} = \op{Sp} (
\Box_{m+1}) + B - (m +1).
\end{gather}
Moreover, since the degree of $L \otimes K^{-m}$ is $(B- m) 2 (\textsl{g} -1)$, by
the Riemann-Roch theorem, $h^0 (L \otimes K^{-m}) := \op{dim} (\ker \Box_m)$ is
positive when $B - m \geqslant 1$, and equal to $(B-m) 2 (\textsl{g}-1) + 1 - \textsl{g}$ when
$B- m \geqslant 2 $.

We deduce the first part of the theorem as follows: by the first identity of
\eqref{eq:bochner_bosonic}, $\Delta = \Box_0 + \frac{1}{2} B$ and if $B
\geqslant 1$, $\la_0 = \tfrac{1}{2} B$ and its multiplicity is $h^0 (
L)$. If $B\geqslant 2$, then $H^{0} (\Sigma, L \otimes K^{-1}) >0$, so
$\la_1 =\frac{1}{2}B + B -1$ by \eqref{eq:rec} and its multiplicity is $H^{0}
(\Sigma, L \otimes K^{-1})$. We~can repeat this until we can no longer apply the Riemann-Roch theorem.

When $L = K^r$, after $r$ iterations, it~comes that $\op{Sp}(\Delta) \setminus \{ \la_0, \ldots, \la_{r-1} \}$
is equal to the spectrum of $\Box_r + \frac{1}{2} r^2$. By Bochner identity,
$\Box_r$ is the Laplace-Beltrami operator.
\end{proof}

\section{Twisted semiclassical calculus}

\label{section:microlocal}

We now introduce the semiclassical limit we will be working with. In~this section, $\Sigma$ need just be a closed manifold.
Let $L \to \Sigma$ be a Hermitian line bundle with a Hermitian connection $\nabla$. We~do not
make any assumption on the curvature $F_{\nabla}$. This section is organized as follows:
\begin{itemize}

\item In Section~\ref{ssection:uniform-pseudo}, we~introduce the twisted semiclassical pseudodifferential calculus;

\item In Section~\ref{ssection:defect-measures}, we~discuss semiclassical defect measures associated to quasimodes in this calculus.
\end{itemize}

\subsection{Definition and main properties}

\label{ssection:uniform-pseudo}

For $m \in
\R$, denote by $\Psi^m_{\mathrm{sc}}(\Sigma)$ the space of standard
($h$-dependent) semiclassical pseudodifferential operators of degree $m \in
\R$ (see \cite[Ch.\,14]{Zworski-12} for an introduction). We~emphasize that
$A \in \Psi^m_{\mathrm{sc}}(\Sigma)$ is a \emph{family} of operators $A =
(A_h)_{h > 0}$, where $A_h = \Op_h(a_h) + \mc{O}(h^\infty)$, $\Op_h$ is an
arbitrary semiclassical quantization on $\Sigma$, and $a_h \in
S^m_h(T^*\Sigma)$ satisfies uniform symbolic estimates in $h > 0$ (see \cite[Ch.\,3]{Lefeuvre-book} for a definition of symbol classes). In~addition, the family $A$ might not be defined for all values of $h > 0$, but only on a subset.

\begin{definition}[Twisted semiclassical quantization]
A family of operators $\bA := (\bA_k)_{k \geq 0}$ such that
\[
\bA_{k} : C^\infty(\Sigma, L^{k}) \to C^\infty(\Sigma, L^{k}),
\]
belongs to the space of \emph{twisted} semiclassical operators $\Psi^m_{\mathrm{tsc}}(\Sigma)$ of degree $m \in \R$ if the following holds: for all contractible open subset $U \subset \Sigma$, for all $\chi,\chi' \in C^\infty_{{\comp}}(U)$, for all $s \in C^\infty(U,L)$ such that $|s|=1$ fiberwise, there exists a family of semiclassical operators $A := (A_h) \in \Psi^m_{\mathrm{sc}}(\Sigma)$ such that for all $f \in C^\infty_{{\comp}}(U)$:
\begin{equation}
\label{equation:def}
s^{-k}\bA_k(f s^{k}) = A_{1/k} f.
\end{equation}
\end{definition}

This quantization was introduced by the first author \cite{Charles-00}. We
also refer to \cite[\S 3.2]{Cekic-Lefeuvre-24} and \cite[\S 2]{Charles--23} for a general introduction to this quantization procedure. We~now recall some of its properties.

Define on the open set $U$ the real-valued $1$-form $\beta \in C^\infty(U,T^*U)$ by $\nabla s = -i \beta \otimes s$ (hence $\nabla s^{k} = -i k \beta \otimes s^{k}$). We~also introduce the symplectic $2$-form
\begin{equation}
\label{equation:omega}
\Omega := dA + i \pi^* F_\nabla,
\end{equation}
where $dA$ is the Liouville $2$-form on $T^*\Sigma$ (see \eqref{equation:liouville}) and $\pi : T^*\Sigma \to \Sigma$ is the projection. (For the magnetic Laplacian on surfaces, we~will have $\Omega = dA + B \pi^*\vol$.) It can be easily checked that \eqref{equation:omega} defines a non-degenerate (closed) $2$-form. Given $p \in C^\infty(T^*\Sigma)$, the Hamiltonian vector field of $p$ associated to $\Omega$ is denoted by $H_p^\Omega$. It is defined through the relation
\[
dp = \Omega(\csbullet, H_p^\Omega).
\]
Finally, the Poisson bracket of two observables $p,q \in C^\infty(T^*M)$ computed with respect to $\Omega$ is written as $\{p,q\}^\Omega =: H^\Omega_p q$. We~sum up the main properties of this calculus (below, $S^m_k(T^*\Sigma)$ denotes the space of symbols of order $m$ which may depend on $k \geq 0$; however, the symbolic estimates are required to be uniform in $k$):

\begin{proposition}
The following properties hold:

\begin{enumerate}

\item \emph{Algebra property.} The space
\[\textstyle
\Psi^\csbullet_{\mathrm{tsc}}(\Sigma) := \bigcup_{m \in \R} \Psi^m_{\mathrm{tsc}}(\Sigma)
\]
is a graded algebra. Namely, for all $\bA\in\Psi^m_{\mathrm{tsc}}(\Sigma)$, $\bB\in\Psi^{m'}_{\mathrm{tsc}}(\Sigma)$, $\bA \circ \bB \in \Psi^{\color{black}m+m'}_{\mathrm{tsc}}(\Sigma)$.

\item \emph{Principal symbol.} For $\bA \in \Psi^m_{\mathrm{tsc}}(\Sigma)$, there is a well-defined principal symbol given by
\[
\sigma_{\bA}(x,\xi ; k) = a(x,\xi + \beta(x)) \in
S^m_{k}(T^*\Sigma)/k^{-1}S^{m-1}_{k}(T^*\Sigma),
\]
where $a$ is the principal symbol of $A_{1/k}$.

\item \emph{Commutator.} For $\bA \in \Psi^m_{\mathrm{tsc}}(\Sigma)$, $\bB \in \Psi^{m'}_{\mathrm{tsc}}(\Sigma)$, $[\bA,\bB] \in k^{-1}\Psi^{m+m'-1}_{\mathrm{tsc}}(\Sigma)$ with principal symbol
\[
\sigma_{k[\bA,\bB]} = -i \{\sigma_{\bA},\sigma_{\bB}\}^\Omega.
\]

\item \emph{Calderon-Vaillancourt.} Let $\bA \in \Psi^0_{\mathrm{tsc}}(\Sigma)$. Then $\bA_k : L^2(\Sigma,L^k) \to L^2(\Sigma,L^k)$ is bounded for all $k \geq 0$ with norm
\[
\|\bA_k\|_{L^2 \to L^2} \leq \|\sigma_{\bA}(\csbullet;k)\|_{L^\infty(T^*\Sigma)} + \mc{O}(k^{-1}).
\]

\end{enumerate}

\end{proposition}

By construction, the operator $\Delta_k := \tfrac{1}{2} {\color{black}k^{-2}} (\nabla^k)^*\nabla^k \in \Psi^2_{\mathrm{tsc}}(\Sigma)$ belongs to this calculus and has principal symbol $p(x,\xi) := \sigma_{\Delta_k}(x,\xi) = \tfrac{1}{2}|\xi|^2_g$.

Conversely, one defines a quantization map $\Op : S^m_{k}(T^*\Sigma) \to \Psi^m_{\mathrm{tsc}}(\Sigma)$ as follows. Consider a smooth function $\chi \in C^\infty(\Sigma \times \Sigma)$, equal to $1$ near the diagonal $\Delta \subset \Sigma \times \Sigma$ and supported in $\{d(x,y) < \iota(g)/2026\}$, where $\iota(g)$ is the injectivity radius of the metric $g$. For $k \geq 0$, and $x,y \in \mathrm{supp}(\chi)$, let $\tau_{x \to y} : L^k_x \to L^k_y$ denote the parallel transport with respect to $\nabla$ along the unique $g$-geodesic joining $x$ to $y$. We~then set:
\[
\Op_{k}(a_{k}) f(x) := \dfrac{1}{2\pi} \int_{\Sigma} \int_{T^*_x\Sigma} e^{-ik\xi(\exp^{-1}_x(y))} a_{k}(x,\xi) \tau_{y \to x}(f(y)) \chi(x,y) \dd\xi \dd y,
\]
where $\dd y$ is the Riemannian volume, and $\dd\xi$ the induced volume in the fibers of $T^*\Sigma$. It is a standard calculation to verify that
\[
\sigma_{\Op_{k}(a_{k})} = [a_{k}] \in S^m_{k}(T^*\Sigma)/k^{-1}S^{m-1}_{k}(T^*\Sigma).
\]

An operator $\bA \in \Psi^\csbullet_{\mathrm{tsc}}(\Sigma)$ is \emph{microlocally compactly supported} if $\bA = \Op_k(a_k) + \mc{O}(k^{-\infty})$ where $a_k \in C^\infty_{\comp}(T^*\Sigma)$ has (uniformly in $k$) compact support on $T^*\Sigma$, and the remainder term $\mc{O}(k^{-\infty})$ denotes an operator with smooth Schwartz kernel on $\Sigma \times \Sigma$, all of whose derivatives are $\mc{O}(k^{-\infty})$ in $L^\infty$-norm. We~denote by $\Psi^{\comp}_{\mathrm{tsc}}(\Sigma)$ the set of microlocally compactly supported operators.

Finally, we~shall use the radial compactification $\overline{T^*\Sigma} :=
T^*\Sigma \sqcup \partial_\infty T^*\Sigma$ of cotangent space (obtained by
adding a sphere at infinity to each fiber of $T^*\Sigma$), see \cite[App.\,E.1.3]{Dyatlov-Zworski-book}. As in standard semiclassical theory, an operator
$\bA \in \Psi^m_{\mathrm{tsc}}(\Sigma)$ is \emph{elliptic} at $(x_0,\xi_0)
\in \overline{T^*\Sigma}$ if there exists a constant $c > 0$ such that for all
$k \geq 0$ large enough, $|\sigma_{\bA}(x,\xi;k)| \geqslant
c\langle\xi\rangle^m$ in a neighborhood of $(x_0, \xi_0)$. The \emph{elliptic set} of an operator is denoted by $\Ell(\bA) \subset \overline{T^*\Sigma}$ and the \emph{characteristic set} is the complement of the elliptic set. Operators can be inverted modulo $\mc{O}(k^{-\infty})$ smoothing operators on their elliptic set (parametrix construction).

\subsection{Defect measures}

\label{ssection:defect-measures}

Quasimodes of energy $E \geq 0$ are sections $u_{k} \in C^\infty(\Sigma,L^k)$ satisfying
\begin{equation}
\label{equation:qm}
(k^{-2}\Delta_k-E) u_{k} = \mc{O}_{L^2}(\eps_{k}), \qquad \|u_k\|^2_{L^2}=1,
\end{equation}
where $\eps_{k} \to_{k \to \infty} 0$. Up to extraction along a subsequence $(k_n)_{n \geq 0}$, there exists a measure $\mu$ on $T^*\Sigma$ such that for all $a \in C^\infty_{{\comp}}(T^*\Sigma)$,
\begin{equation}
\label{equation:limit}
\lim_{n \to \infty} \langle\Op_{k_n}(a)u_{k_n},u_{k_n}\rangle_{L^2} = \int_{T^*\Sigma} a(x,\xi) d\mu(x,\xi).
\end{equation}
That $\Delta_k$ is elliptic on $\partial_\infty T^*\Sigma$ (boundary at
infinity of the radial compactification of~$T^*\Sigma$), implies that $\mu$ is
a probability measure, see \cite[App.\,E, Exer.\,23]{Dyatlov-Zworski-book}. In~the following, recall that $p(x,\xi) = \tfrac{1}{2}|\xi|^2_g$.

\begin{proposition}
\label{proposition:scl-measure}
Let $\mu$ be a probability measure associated to the quasimodes \eqref{equation:qm}. Then the following holds:
\begin{enumerate}

\item\label{proposition:scl-measurei} Assume that $\eps_k = o(1)$. Then
\[
\supp(\mu) \subset \mathrm{Char}(k^{-2}\Delta_k-E) = \{p=E\}.
\]

\item\label{proposition:scl-measureii} Assume that $\eps_k = o(1/k)$. Then $\mu$ is invariant by the magnetic Hamiltonian vector field $H_{p}^\Omega$, that is for all $a \in C^\infty_{{\comp}}(T^*\Sigma)$,
\[
\int_{T^*\Sigma} H_{p}^\Omega a(x,\xi) ~\dd\mu(x,\xi) = 0.
\]
\end{enumerate}
\end{proposition}

We refer to \cite[App.\,E.3]{Dyatlov-Zworski-book} for a proof.

\subsection{Technical estimates}

We will need some precise estimates on the remainder in the standard operations of semiclassical analysis (product, commutator, propagation, etc.). Recall that $(\Phi_t)_{t \in \R}$ is the magnetic flow on $T^*\Sigma$, that is the Hamiltonian flow generated by $H_p^\Omega$.

\begin{proposition} For any $E>0$, there exists a constant $C > 0$ such that for all $k \geq 0$, for all $a,b \in
C^\infty_{\comp}(T^*\Sigma)$ with $\mathrm{supp}(a), \mathrm{supp}(b) \subset
\{|\xi|_g \leq E\}$, the following holds:
\begin{align}
\label{equation:calderon-fine}
\|\Op_k(a)\|_{L^2 \to L^2} & \leq \|a\|_{L^\infty(T^*\Sigma)} + k^{-1/2}\|a\|_{C^{14}(T^*\Sigma)}, \\
\label{equation:product-fine}
\|\Op_k(a)\Op_k(b)-\Op_k(ab)\|_{L^2 \to L^2} &\leq Ck^{-1} \|a\|_{C^{17}(T^*\Sigma)}\|b\|_{C^{17}(T^*\Sigma)}, \\
\label{equation:commutator-fine}
\|[\Op_k(a),\Op_k(b)]-\Op_k(-ik^{-1}\{a,b\})\|_{L^2 \to L^2} &\leq Ck^{-1}\|a\|_{C^{17}(T^*\Sigma)}\|b\|_{C^{17}(T^*\Sigma)}, \\
\label{equation:egorov-fine}
\|e^{it k^{-1} \Delta} \Op_k(a) e^{-itk^{-1}\Delta}-\Op_k(a \circ \Phi_t)\|_{L^2 \to L^2} &\leq C k^{-1} \int_0^t \|a \circ \Phi_s\|_{C^{17}(T^*\Sigma)} \dd s.
\end{align}

\end{proposition}

The above estimates are very likely suboptimal by far; however, we~did not try to optimize them.

\begin{proof}
Estimate \eqref{equation:calderon-fine} follows from \cite[Th.\,5.26]{Nonnenmacher-notes} applied in dimension $2$ and Lemma \cite[Lem.\,A.4]{Dyatlov-Jin-Nonnenmacher-22}. The estimates \eqref{equation:product-fine} and \eqref{equation:commutator-fine} can be found in \cite[Lem.\,A.6]{Dyatlov-Jin-Nonnenmacher-22} for the standard semiclassical quantization. It is an exercise to verify that the proofs go through for the twisted quantization. Finally, \eqref{equation:egorov-fine} follows from \eqref{equation:commutator-fine}. Indeed, since $e^{it k^{-1} \Delta}$ is unitary on $L^2(\Sigma, L^k)$, it~suffices to estimate $\|B(t)\|$, where
\[
B(t) := \Op_k(a) - e^{-itk^{-1}\Delta}\Op_k(a\circ\Phi_t)e^{itk^{-1}\Delta}.
\]
Notice that $B(0)=0$ and using \eqref{equation:commutator-fine}:
\[
\begin{split}
B'(t) &= e^{-itk^{-1}\Delta}\left(k[ik^{-2}\Delta_k,\Op_k(a\circ\Phi_t)]-\Op_k(H_p a \circ\Phi_t) \right) e^{itk^{-1}\Delta} \\
& = \mc{O}_{L^2 \to L^2}(\|a\circ\Phi_t\|_{C^{17}(T^*\Sigma)} h),
\end{split}
\]
where $p$ is the principal symbol of $k^{-2}\Delta_k$. Writing
\[
B(t) = \int_0^t B'(s)\dd s,
\]
we find the claimed result.
\end{proof}

\section{Construction of eigenstates}

\label{section:constant}

We now assume that $(\Sigma,g)$ is a Riemannian surface of constant curvature $-1$, genus $\geq 2$, and $L \to \Sigma$ is a Hermitian line bundle equipped with a unitary connection $\nabla$ with curvature $F_\nabla = -iB\vol$ and $B > 0$ is constant. Without loss of generality, since $H^2(\Sigma,\Z) \simeq \Z$ and we will be
considering power $L^k \to \Sigma$, we~can further assume that $\deg(L)=1$.
Note that by \eqref{equation:deg}, this forces $B = |\chi(\Sigma)|^{-1}$. Hence, in what follows, the magnetic intensity on $M$ \emph{is fixed and equal to $B = |\chi(\Sigma)|^{-1}$} on $L \to \Sigma$.

This section is organized as follows:
\begin{itemize}

\item In Section~\ref{ssection:averaging}, we~adapt Weinstein's averaging method \cite{Weinstein-77} to our context and introduce an operator $\bA_k$ whose principal symbol generates a $2\pi$-periodic flow on $\{p < E_c\}$ which is a renormalization of the magnetic flow;

\item We then use it in Section~\ref{ssection:quasimodes} to construct specific eigenstates concentrating in phase space on any closed orbit of the magnetic flow.
\end{itemize}

\subsection{The averaging method} \label{ssection:averaging} By Proposition \ref{theo:Landau_Level}
applied to $L^k \to \Sigma$ with curvature $F_{\nabla^k}=-ikB
\op{vol}$, the first eigenvalues of $\Delta_k = \tfrac{1}{2}\nabla^*\nabla : C^\infty(\Sigma,L^k) \to C^\infty(\Sigma,L^k)$ are given for $m < \lfloor kB\rfloor $ by
\begin{equation}
\label{equation:ev}
\lambda_{k,m} = kB (m+\sfrac{1}{2})-\sfrac{m(m+1)}{2}.
\end{equation}
Notice that $\lambda_{k, \lfloor kB \rfloor - 1 } \sim k^2 E_c$ as $k
\rightarrow \infty$, where $E_c
= \tfrac{1}{2} B^2$ is the critical energy. As a consequence, for any $\eps >0$, we~have that
\[
\op{Sp}
(k^{-2} \Delta_k) \cap [0, E_c -\eps ] \subset \{k^{-2} \la_{k,m} ~|~ m < \lfloor
kB\rfloor \}\]
when $k$ is sufficiently large.

\subsubsection{Weinstein's periodic operator} The following paragraph is
inspired by Weinstein's averaging method \cite{Weinstein-77}. Let $\Pi_{k,m}$ be the orthogonal projector of $C^\infty(M, L^k)$ onto
$\ker(\Delta _k -\lambda_{k,m})$, and set:
\begin{equation}
\label{equation:a}
\bA_k := k^{-1} \sum_{m =0 }^{ \lfloor kB \rfloor -1} m \Pi_{k,m}.
\end{equation}
Observe that by construction $\op{Sp}(\bA_k) \subset k^{-1} \Z_{\geq 0}$ and thus
\begin{equation}
\label{equation:cool}
e^{2i k \pi \bA_k } = \mathbf{1}.
\end{equation}
Furthermore, by~\eqref{equation:ev}, on the space
\[
\mc{I}_k := \bigoplus_{m=0}^{ \lfloor kB \rfloor -1}
\ker(\Delta_k-\lambda_{k,m}),
\]
the operator $\bA_k$ satisfies the identity
\begin{equation}
\label{equation:relation}
k^{-2} \Delta_k = B(\bA_k+ \tfrac{1}{2} k^{-1})- \tfrac{1}{2}
\bA_k(\bA_k+ k^{-1}).
\end{equation}
Equivalently,
\begin{equation}
\label{equation:relation2}
\bA_k= B-\tfrac{1}{2} k^{-1} - \sqrt{ B^2 - 2 k^{-2}
\Delta_k + \tfrac{1}{4} k^{-2}},
\end{equation}
on $\mc{I}_k$. Finally, observe that
\begin{equation}
\label{eq:integrale_projecteur}
\Pi_{k,m} = \dfrac{1}{2\pi} \int_0^{2\pi} e^{-imt} e^{itk\bA_k} \dd t.
\end{equation}
This turns out to be a Fourier Integral Operator in the semiclassical regime $k \to \infty$.

Let
\[
D^*\Sigma := \{(x,\xi) \in T^*\Sigma ~|~ |\xi| < B\}.
\]
Let $a$ be the function defined on $D^*\Sigma$ through the relation
\begin{equation}
\label{equation:mercredi-aprem}
p(x,\xi) = \beta (a (x, \xi)), \qquad \forall (x,\xi) \in D^*\Sigma,
\end{equation}
where $\beta : [0,B] \rightarrow [0, \frac{1}{2} B^2 ]$ is the function $ \beta (s) := B s -\frac{1}{2} s^2$. The following holds:

\begin{lemma}
Let $\chi \in C^\infty_{{\comp}}(T^*\Sigma)$ such that $\supp(\chi) \subset D^*\Sigma$ and $\varphi \in C^\infty(\R)$ such that $\supp(\varphi) \subset (-\infty,B)$. Then
\[
\Op_k(\chi) \bA_k, \varphi(\bA_k) \in \Psi^{\comp}_{\mathrm{tsc}}(\Sigma),
\]
and these operator have respective principal symbols $\chi a$ and $\varphi(a)$.
\end{lemma}

We emphasize that $\bA_k$ is \emph{not} a pseudodifferential operator because of the spectral cutoff involved in its definition (the spectral projector onto $\mc{I}_k$). However, in what follows, we~will always use $\Op_k(\chi) \bA_k, \varphi(\bA_k)$ which \emph{are} pseudodifferential operators. For simplicity, we~also say that $a$ is the principal symbol of $\bA_k$, although this only makes sense in $D^*\Sigma$.

\begin{proof}
We define $f_k : [0,\tfrac{1}{2}B^2] \to \R$ by:
\begin{equation}
\label{equation:fk}
f_k(s) = B-\tfrac{1}{2}k^{-1}-\sqrt{B^2-2s+\tfrac{1}{4}k^{-2}}.
\end{equation}
Notice that $\varphi \circ f_k$ is smooth with compact support in $[0,\tfrac{1}{2}B^2)$ by the assumption made on the support of $\varphi$; it can thus be extended to a smooth compactly supported function $c_k$ on $[0,\infty)$. In~addition, using \eqref{equation:relation2}, the equality $\varphi(\bA_k) =\varphi(f_k(k^{-2}\Delta_k)) = c_k(k^{-2}\Delta_k)$ holds. The function $c_k$ has compact support on $[0,\infty)$ and uniformly bounded derivatives as $k \to \infty$. Hence, by~standard functional calculus (see \cite[Th.\,14.9]{Zworski-12} for instance), $c_k(k^{-2}\Delta_k) = \varphi(\bA_k) \in \Psi^{\comp}_{\mathrm{tsc}}(\Sigma)$.

We now establish that $\Op_k(\chi) \bA_k \in \Psi^{\comp}_{\mathrm{tsc}}(\Sigma)$. For that, we~write
\[
\Op_k(\chi) \bA_k =\Op_k(\chi) \psi(k^{-2}\Delta_k) \bA_k + \Op_k(\chi) (\mathbbm{1}-\psi(k^{-2}\Delta_k)) \bA_k,
\]
where $\psi $ is smooth with compact support in $(-\infty, E_c)$ and such that $\psi(p) \equiv 1$ on the support of $\chi$. Notice that $\Op_k(\chi) \psi(k^{-2}\Delta_k) \bA_k = \Op_k(\chi) d_k(k^{-2}\Delta_k)$, where $d_k(x) = \psi(x) f_k(x)$ and $f_k$ was defined in \eqref{equation:fk}. The function $d_k$ extends to a smooth compactly supported function on $[0,\infty)$ and the same argument as in the previous paragraph using functional calculus shows that $\Op_k(\chi) \psi(k^{-2}\Delta_k) \bA_k \in \Psi^{\comp}_{\mathrm{tsc}}(\Sigma)$. Let us now prove that $\Op_k(\chi) (\mathbbm{1}-\psi(k^{-2}\Delta_k)) \bA_k$ is a residual operator, namely that it has smooth Schwartz kernel all of whose derivatives are $\mc{O}(k^{-\infty})$. By standard arguments, it~suffices to prove that
\[
\Op_k(\chi) (\mathbbm{1}-\psi(k^{-2}\Delta_k)) \bA_k : H^{-N}_k(\Sigma,L^k) \to H^{+N}_k(\Sigma,L^k)
\]
is bounded for all $N \geq 0$ with norm $\mc{O}(k^{-\infty})$, where the space $H^s(\Sigma,L^k)$ is defined as the completion of $C^\infty(\Sigma,L^k)$ with respect to the norm
\begin{equation}
\label{equation:sobolev-norm}
\|u\|^2_{H^s(\Sigma,L^k)} := \|(k^{-2}\Delta_k+\mathbbm{1})^{s/2} u\|^2_{L^2(\Sigma,L^k)}.
\end{equation}
By the support property of $\psi$, the operator $\Op_k(\chi) (\mathbbm{1}-\psi(k^{-2}\Delta_k))$ clearly satisfies this property (this is a composition of two operators with disjoint microsupport). Hence, it~suffice to show that $\bA_k : H^{-N}_k(\Sigma,L^k) \to H^{-N}_k(\Sigma,L^k)$ is bounded with norm $\mc{O}(1)$ {\color{black}for all integers $N \geq 0$}. Using \eqref{equation:sobolev-norm}, this amounts to proving that
\[
(k^{-2}\Delta_k+\mathbbm{1})^{-N/2} \bA_k (k^{-2}\Delta_k+\mathbbm{1})^{N/2} : L^2(\Sigma,L^k) \to L^2(\Sigma,L^k)
\]
has norm $\mc{O}(1)$. However, this is immediate using \eqref{equation:a} and
the fact that this operator acts diagonally on the Fourier decomposition into eigenmodes of $\Delta_k$.

Finally, to compute the principal symbol of these operators, merely observe
that by \eqref{equation:relation}, we~have $p = B a - a^2/2 = \beta(a)$.
\end{proof}

The function $a$ will play an important role in the sequel. One important property is the following:

\begin{lemma}
\label{lemma:2pi-periodic}
The Hamiltonian vector field $H_a^\Omega$ of the symbol $a$ generates a $2 \pi$\nobreakdash-peri\-odic flow on $D^*\Sigma$.
\end{lemma}

\begin{proof}This can be seen by writing $a = \al (p)$ with $\al(y) := B - \sqrt{ B^2 - 2 y }$, and noticing
that $\al' (E) = (B^2 - 2 E)^{-\sfrac{1}{2}}$ is the period of the orbits of $H_p^\Omega$
with energy $p =E$ (see Proposition \ref{proposition:dynamic}, item \eqref{proposition:dynamici}).
\end{proof}

Given a function $f \in C^\infty (D^*\Sigma)$, introduce its average
$\langle f \rangle \in C^\infty (D^*\Sigma)$ by:
\begin{equation}
\label{equation:average}
\langle f \rangle (z) := \frac{1}{2\pi} \int_0^{2 \pi} f (\Phi_t^a (z)) ~\dd t,
\end{equation}
where $(\Phi_t^a)_{t \in \R}$ is the $2\pi$-periodic Hamiltonian flow generated by $H_a^\Omega$ on $D^*\Sigma$:

\begin{lemma} \label{lem:commutateur_espace_propre}
Fix $\eps > 0$. For all $f \in C^\infty_{{\comp}} (D^*\Sigma)$, $k \geq 0$, for all $0 \leq m \leq (1-\eps)kB$, the following holds:
\[
\Pi_{k,m}\Op_{k}(f) \Pi_{k,m} = \Op_{k}(\langle f \rangle)
\Pi_{k,m} + \bigo_{\mc{L}(L^2)}(k^{-1}).
\]
Here the $\bigo $ depends on $f$ but is uniform with respect to $m$.
\end{lemma}

The proof is a variation on Egorov's theorem (see \cite[Th.\,15.2]{Zworski-12} for instance).

\begin{proof}
In the proof, the $\mc{O}$'s are measured in operator norm $L^2 \to L^2$. Fix
$f \in C^\infty_{{\comp}}(D^*\Sigma)$ and choose $\varphi \in
C^\infty(-\infty, B)$
such that $ \varphi (a) =1$ on a neighbourhood of the support of $f$. Then
\[
\Op_k(f) \varphi (\bA_k) \equiv \Op_k(f) \equiv \varphi (\bA_k) \Op_k
(f),
\]
where $\equiv$ stands for equality modulo $\bigo (k^{-\infty})$. This yields
\begin{xalignat*}{2}
[ \Op_k(f), \bA_k ] & \equiv [ \Op_k(f), \varphi (\bA_k) \bA_k ]
= \frac{1}{ik} \Op_k (\{ f, \varphi (a) a \}) + \bigo (k^{-2}) \\ & =
\frac{1}{ik} \Op_k (\{ f, a \}) + \bigo (k^{-2}).
\end{xalignat*}
Write $U_t := e^{itk\bA_k}$. Integrating the previous relation, we~obtain the weak form of Egorov's theorem for $\bA_k$:
\begin{gather} \label{eq:egorov}
U_t \Op_k(f) U_{-t} = \Op_k(f \circ \Phi_t^a) + \bigo (k^{-1}).
\end{gather}
Since $U_t \Pi_{k,m} = \Pi_{k,m} U_t = e^{itm} \Pi_{k,m}$, we~obtain that
\[
\Pi_{k,m} \Op_k(f) \Pi_{k,m} = \Pi_{k,m} U_t \Op_k(f) U_{-t} \Pi_{k,m} = \Pi_{k,m} \Op_k(f \circ
\Phi_t^a) \Pi_{k,m} + \bigo (k^{-1}).
\]
Averaging over time, we~find that
\begin{gather}
\label{eq:moyenne}
\Pi_{k,m} \Op_k(f) \Pi_{k,m} = \Pi_{k,m} \Op_k(\langle f \rangle) \Pi_{k,m} + \bigo (k^{-1}).
\end{gather}
Moreover, since $\langle f \rangle \circ \Phi_t^a = \langle f \rangle $, we
deduce from \eqref{eq:egorov} that $ [ U_t, \Op_k(\langle f \rangle) ] =
\bigo (k^{-1})$ and from \eqref{eq:integrale_projecteur} that $[\Pi_{k,m}, \Op_k(
\langle f \rangle)] = \bigo (k^{-1})$. The claim then follows from this last
relation and \eqref{eq:moyenne}.
\end{proof}
\subsection{Localized eigenstates} \label{ssection:quasimodes} We now construct specific eigenstates for the operator $k^{-2}\Delta_k$ with adapted localization properties.

\subsubsection{Energy $E=0$} We first investigate the low energy eigenstates in the semiclassical regime
\[
k^{-2}\Delta_k u_k = E_k u_k, \qquad u_k \in C^\infty(\Sigma,L^k),\; \|u_k\|_{L^2}=1,\; E_k = \mc{O}(1/k).
\]
This corresponds the regime of eigenvalues $\lambda_{k,m}$ with $k \to \infty$ and $m \leqslant C$, where $C > 0$ is a fixed constant. From now on, we~use the notation $h = 1/k$. Let $\underline{x} \in \Si$, $u \in L_{\underline{x}}$ with $|u|=1$ and define the Dirac section of
$L^k$ by $(\delta_u^k, s) = s(\underline{x})/u^k$ for any $k \geq 0$, $s \in C^\infty (\Sigma,L^k)$. We~introduce the section
\begin{gather} \label{eq:def_etat_coherent}
\bme_{k,m,\underline{x}} := \Pi_{k,m} \delta_u^k.
\end{gather}

The following holds:

\begin{proposition}
\label{proposition:e0}
Let $C > 0$ and $f \in C^\infty_{{\comp}}(T^* \Sigma)$. Then for all $k \geq 0$, $0 \leq m \leq C$, $\bme_{k,m,\underline{x}} \in C^\infty(\Sigma,L^k)$ satisfies
\begin{xalignat}{2}
\begin{split} \label{eq:estim1}
\| \bme_{k,m,\underline{x}} \|^2 &= \frac{k}{2\pi} (1 + \bigo (k^{-1/2})), \\
\langle \Op_k(f) \bme_{k,m,\underline{x}}, \bme_{k,m,\underline{x}} \rangle &= \| \bme_{k,m,\underline{x}} \|^2 \bigl(f (\underline{x},0) + \bigo
(k^{-1/2}) \bigr).
\end{split}
\end{xalignat}
\end{proposition}

As we will see, $\bme_{k,m,\underline{x}}$ is a Gaussian state centered at
$\underline{x}$. We~provide a proof building on \cite{Landau1,Landau2}.

\begin{proof}
The section $ \bme_{k,m,\underline{x}}$ is given in terms of the Schwartz kernel of
the projector $\Pi_{k,m}$ by
\begin{gather} \label{eq:noyau_coherent}
\bme_{k,m,\underline{x}} (x) \otimes \con{u}^k = \Pi_{m,k} (x, \underline{x})
\in L_x^k \otimes \con{L}_{\underline{x}}^k.
\end{gather}
It was proved in \cite{Landau1,Landau2} that $\Pi_{k,m} (x,\underline{x})$ has an
asymptotic expansion, which leads to the following expansion for $ \bme_{k,m,\underline{x}}$. Introduce coordinates $(x_1, x_2)$ centered at $\underline{x}$ and such that the metric $g$
is $(1 + \bigo (|x|)) (dx_1^2 + dx_2^2)$ with $|x|^2 = x_1^2 + x_2^2$. Choose a primitive $\al = \al_1
dx_1 + \al_2 dx_2 $ of the Riemannian volume $\vol$ on a neighborhood of $y$, such that $ \al_1 (\underline{x})
= \al_2 (\underline{x}) =0$ and $\partial_{x_i} \al_j (\underline{x})$ is antisymmetric. Let $t$ be a
frame of $L$ such that $t (\underline{x}) = u$ and $\nabla t = \tfrac{1}{i}
\al \otimes t$. Then one has on a neighborhood of~$\underline{x}$:
\begin{equation} \label{eq:theta}
\bme_{k,m,\underline{x}} (x) = \frac{k}{2 \pi} e^{-\frac{1}{4} k |x|^2 } \sum_{\ell
=0}^N k^{ -\sfrac{\ell}{2} } a_\ell (k^{\sfrac{1}{2} } x_1, k^{\sfrac{1}{2}}
x_2) t^k(x) + \bigo (k^{- \psfrac{N+1}{2}}),
\end{equation}
where the $\bigo $ is in the $\mathcal{C}^0$-norm and the $a_{\ell}$'s are polynomials depending on $m$ and smoothly on
$\underline{x}$. The leading coefficient $a_0$ is given in terms of the $m$-th
Laguerre polynomial $Q_m$ by
\begin{gather} \label{eq:leading}
a_0 (x_1, x_2) = Q_m (|x|^2), \qquad Q_m (t) = \frac{1}{m!} \Bigl(\frac{d}{dt}-1\Bigr)^m t^m.
\end{gather}
Moreover, $\bme_{k,m,\underline{x}}$ is $\bigo_{C^\infty} (k^{-\infty})$ on any compact set not
containing $\underline{x}$.
Formulas \eqref{eq:theta}, \eqref{eq:leading} correspond to \cite[Eqs.\,(26) \& (40)]{Landau1}.

Since $\Pi_{k,m}$
is a projector, $\| \bme_{k,m,\underline{x}} \|^2 = \Pi_{k,m}
(\underline{x},\underline{x})$, so~by \eqref{eq:theta},
\[
\| \bme_{k,m,\underline{x}} \|^2 = \frac{k}{2\pi} + \bigo (k^{1/2}),
\]
because $Q_m (0) =1$. We~then compute the asymptotic expansion of the scalar product
\[
\langle \Op_k(f) \bme_{k,m,\underline{x}},
\bme_{k,m,\underline{x}} \rangle = \sum_{\ell, \ell' \geqslant 0 }
k^{-\frac{1}{2} (\ell + \ell')} I_{\ell, \ell'}(f,k)
\]
with $I_{\ell, \ell'} (f,k)$ given by the
integral
\[
\Bigl(\frac{k}{2 \pi} \Bigr)^4 \!\!\int\! e^{ik (x- y) \xi - \frac{1}{4}k (|x|^2 + |y^2|) }
f (\tfrac{1}{2} (x+y), \xi) a_{\ell} (k^{\sfrac{1}{2}} x) \overline{
a_{\ell'} (k^{\sfrac{1}{2}} y)} \; \dd\xi_1 \dd\xi_2 \dd x_1 \dd x_2 \dd y_1 \dd y_2.
\]
Notice that the phase is quadratic, non degenerate. It follows from a stationary phase computation that $I_{\ell, \ell'} (f,k) = k C_{\ell, \ell'} f(0, 0) + \bigo
(k^{\sfrac{1}{2}})$ for some complex coefficient $C_{\ell, \ell'}$. This yields:
\[
\langle \Op_k(f) \bme_{k,m,\underline{x}}, \bme_{k,m,\underline{x}} \rangle
= k f(0,0) C_{0,0} + \bigo (k^{1/2})\]
and $C_{0,0} = (2 \pi)^{-1} $ because of the previous estimate of $\| \bme_{k,m,\underline{x}}
\|^2$. This concludes the proof. We~could actually slightly improve the result
by replacing the $\bigo (k^{-1/2})$'s in~\eqref{eq:estim1} by $\bigo (k^{-1})$'s, which follows
from the fact that
each $a_\ell$ has the same parity as~$\ell$.
\end{proof}

\subsubsection{Energies $0 < E < E_c$} We now investigate the eigenstates such that
\[
k^{-2}\Delta_{k} u_k = (E+o(1))u_k,
\]
with $0 < E < E_c$. Equivalently, this corresponds to the regime of
eigenvalues $\lambda_{k,m}$ with $k \to \infty$ and $\eps k B \leqslant m \leqslant (1- \eps)kB$, where $\eps> 0$ is fixed. {\color{black}Below, recall that $a$ is the principal symbol of $\bA_k$, see \eqref{equation:mercredi-aprem}.}

\begin{proposition} \label{lem:gaussian-beam}
Let $z_0 \in D^*_0\Sigma := D^*\Sigma\setminus \{ \xi = 0 \}$. Then, there exists a neighborhood $U \subset D^*_0\Sigma$ of $z_0$ such that for all $f \in C^\infty_{{\comp}}(T^*\Sigma)$, for all $k \geq 0$, $\eps kB \leq m \leq (1-\eps)kB$, and $z \in U$ such that
\[
m=ka(z),
\]
there exists $\bme_{k,m,z} \in C^\infty(\Sigma,L^{k})$ such that $\| \bme_{k,m,z} \|^2 = 1$ and the
projection
\[
\bmf_{k,m,z} := \Pi_{k,m} \bme_{k,m,z}
\]
satisfies
\begin{xalignat}{2}
\begin{split}
\| \bmf_{k,m,z} \|^2 &= \tfrac{1}{ 2 \sqrt \pi} k^{-\sfrac{1}{2}}(1+o_{f,U}(1)), \\
\langle \Op_k(f) \bmf_{k,m,z}, \bmf_{k,m,z} \rangle &= \|\bmf_{k,m,z}\|^2 \bigl(\langle f \rangle (z) + \bigo_{f,U}
(k^{-\sfrac{1}{4}}) \bigr).
\end{split}
\end{xalignat}
The remainder terms $\bigo$'s are uniform with respect to $k \geq 0$.
\end{proposition}

The state $\bme_{k,m,z}$ is a coherent state centered at $z$ and $\bmf_{k,m,z}$
is a Gaussian beam supported on the magnetic geodesic of $z$ (i.e. the flow line of $H_a^\Omega$ passing through $z$).

\begin{proof}
We use $h=1/k$ in the proof once again. Let $\bme_{k,m,z} \in C^\infty(\Sigma,L^k)$ be a Gaussian state centered at $z$ and such that $\|\bme_{k,m,z}\|_{L^2}=1$ (see \cite[Ex.\,1, Ch.\,5]{Zworski-12} for instance). A~quick computation using the stationary phase lemma reveals
\begin{gather} \label{eq:etat_coherent}
\Op_k(g) \bme_{k,m,z} = g (z) \bme_{k,m,z} + \bigo_{L^2} (
h^{\sfrac{1}{2}})
\end{gather}
for any $g \in C^\infty_{{\comp}}(T^*\Sigma)$. Then $\bmf_{k,m,z} := \Pi_{k,m} \bme_{k,m,z}$ (with $m=ka(z)$) satisfies by Lemma
\ref{lem:commutateur_espace_propre} that
\begin{xalignat}{2} \notag
\langle \Op_k(f) \bmf_{k,m,z}, \bmf_{k,m,z} \rangle & = \langle \Op_k(\langle f \rangle)
\bme_{k,m,z}, \Pi_{k,m} \bme_{k,m,z} \rangle + \bigo (h \|\bmf_{k,m,z}\|) \\
& = \langle f \rangle (z) \|\bmf_{k,m,z}\|^2 + \bigo (h^{\sfrac{1}{2}} \|\bmf_{k,m,z}\|).
\end{xalignat}
To estimate the norm of $\bmf_{k,m,z}$, we~use
\eqref{eq:integrale_projecteur}, namely (recall that $U_t = e^{itk\bA_k}$):
\begin{gather} \label{eq:int_norme}
\begin{split}
\|\bmf_{k,m,z}\|^2 = \langle \Pi_{k,m} \bme_{k,m,z}, \bme_{k,m,z} \rangle & = \frac{1}{2 \pi }
\int_0^{2 \pi } e^{-itm } \langle U_t \bme_{k,m,z},\bme_{k,m,z}
\rangle \; \dd t \\
& = \frac{1}{2 \pi } \int_0^{2 \pi } \langle e^{itk(\bA_k-a(z))}
\bme_{k,m,z},\bme_{k,m,z} \rangle \; \dd t,
\end{split}
\end{gather}
where we used that $m = k a(z)$.
By assumption, the microsupport of $\bme_{k,m,z}$ is $\{z\}$ (that is for
any $\varphi \in C^\infty _{{\comp}} (T^* \Si)$ identically equal to 1 on a
neighborhood of $z$, $\Op_k(\varphi) \bme_{k,m,z} = \bme_{k,m,z} + \bigo_{C^\infty} (
h^{\infty})$). Then by Egorov's theorem, the microsupport of
$U_t \bme_{k,m,z}$ is $ \{ \Phi_t^a(z) \}$. Since $\Phi_t^a (z) = z $ if and only if $t$
is a multiple of $2 \pi$, the integral
\eqref{eq:int_norme} is modified by a $\bigo (h^{\infty})$ if we
restrict the integral to a neighborhood of $t=0$.

Microlocally, $\bA_k-a(z)$ can be conjugated by a
Fourier Integral Operator to $-ih\partial_{x_1}$ acting on
$\R^2_{x_1, x_2}$, see \cite[Th.\,12.3]{Zworski-12}. Since we can restrict the integral \eqref{eq:int_norme} to a
neighborhood of the origin, it~is sufficient to do the computation in this
model. The propagator is then given by $e^{itk(\bA_k-a(z))} \Psi (x_1, x_2) = \Psi (x_1 +t, x_2)$ for $\Psi \in C^\infty(\R^2)$. The point $z$ is in the characteristic of $\bA_k-a(z)$; it is therefore mapped to \hbox{$(\underline{x}, \underline{\xi}_1=0,\underline{\xi}_2) \in T^*\R^2$}. Up to conjugating again by a translation, we~can further assume that $\underline{x}=0$. Let
$\Psi_{\underline{x}, \underline{\xi}} \in C^\infty (\R^2)$ be
the coherent state
\[
\Psi_{\underline{x}, \underline{\xi}} (x) = (\pi h)^{-1/2} e^{ -
\frac{1}{2h} | x - \underline{x}|^2 + \frac{i}{h} \underline{\xi} (x - \underline{x}) } = (\pi h)^{1/2}e^{ -
\frac{1}{2h} | x|^2 + \frac{i}{h} \underline{\xi}_2 x_2 }.\]
By an exact computation, we~find that for $\delta > 0$, $\chi \in C^\infty_{{\comp}}(-\delta,\delta)$ such that $\chi \equiv 1$ in a neighborhood of $0$,
\[
\begin{split} \frac{1}{2 \pi } \int_0^{2 \pi } \langle e^{itk(\bA_k-a(z))} \bme_{k,m,z}&,\bme_{k,m,z} \rangle_{L^2} \chi(t)  \dd t \\
& = \frac{1}{2\pi} \int_{-\infty}^{+\infty} \int_{\R^2} \Psi _{ \underline{x},
\underline{\xi} } (x_1 + t, x_2) \con { \Psi_{ \underline{x}, \underline{\xi} } (x_1,
x_2)} \chi(t) \; \dd x_1 \dd x_2 \dd t \\
& = \frac{ h^{1/2}}{2 \sqrt \pi }\,(1+ \mc{O}(h^\infty)).
\end{split}
\]
Going back to \eqref{eq:int_norme}, we~find that $ \|\bmf_{k,m,z}\|^2_{L^2} =
\tfrac{1}{ 2 \sqrt \pi}\, h^{1/2} + \bigo (h^{\infty})$, which
concludes the proof.
\end{proof}

\section{Proof of the main results}

\label{section:proofs}

In this final section, we~prove Theorems \ref{theorem:main}, \href{https://jep.centre-mersenne.org/item/10.5802/jep.333.pdf#page=6}{1.3} and conclude with the proof of Theorem \href{https://jep.centre-mersenne.org/item/10.5802/jep.333.pdf#page=5}{1.2}. Up to a preliminary rescaling, we~can always assume that $B=1/|\chi(\Sigma)|$.

\subsection{Proof of Theorem \ref{theorem:main}} We can now characterize the semiclassical defect measures in the three regimes.

\skpt
\begin{proof}[Proof of Theorem \ref{theorem:main}]

\eqref{theorem:maini}
Suppose $E=0$. We~claim that any probability measure $\mu$ on $\Sigma$ is a semiclassical limit of eigenstates \eqref{equation:eigenstates}. First, given $x \in \Sigma$,
\[
\bme'_{k,m,x} := \bme_{k,m,x}/\|\bme_{k,m,x}\|_{L^2}
\]
converges semiclassically to $\delta_x$ by Proposition \ref{proposition:e0}. Given $x_1,\dots, x_N \in \Sigma$, $\bE_{k,m,N} := N^{-1/2}(\bme'_{k,m,x_1} +\cdots+ \bme'_{k,m,x_N})$ converges semiclassically to $N^{-1}(\delta_{x_1}+\cdots+\delta_{x_N})$ and $\|\bE_{k,m,N}\|^2_{L^2} = 1 + \mc{O}(k^{-\infty})$. Indeed, to compute the norm, let $\varphi_{i} \in C^\infty(\Sigma)$ be a small bump function centered at $x_i$, equal to $1$ on a small neighborhood of $x_i$ and such that $x_j \notin \supp(\varphi_i)$ for $j \neq i$. It is then immediate to verify that $\varphi_i \bme'_{k,m,x_i} = \bme'_{k,m,x_i} + \mc{O}_{C^\infty}(h^\infty)$ and $\varphi_j \bme'_{k,m,x_i} = \mc{O}_{C^\infty}(h^\infty)$. This leads to
\[
\begin{split}
\|\bE_{k,m,N}\|^2_{L^2} & = N^{-1}\langle \bme'_{k,m,x_1} +\cdots+ \bme'_{k,m,x_N}, \bme'_{k,m,x_1} +\cdots+ \bme'_{k,m,x_N}\rangle_{L^2} \\
& = 1 + N^{-1} \sum_{i \neq j} \langle \bme'_{k,m,x_i},\bme'_{k,m,x_j}\rangle_{L^2} \\
& = 1 + N^{-1} \sum_{i \neq j} \langle \varphi_i \bme'_{k,m,x_i},\bme'_{k,m,x_j}\rangle_{L^2} + \mc{O}(h^\infty) \\
& = 1 + N^{-1} \sum_{i \neq j} \langle \bme'_{k,m,x_i},\varphi_i \bme'_{k,m,x_j}\rangle_{L^2} + \mc{O}(h^\infty) = 1 + \mc{O}(h^\infty).
\end{split}
\]
The same computation shows that the associated defect measure is equal to \hbox{$N^{-1}(\delta_{x_1}+\cdots+\delta_{x_N})$}. Finally, since Dirac masses are dense in probability measures, a diagonal extraction allows to construct a family $(\bE_{k_n,m_n,N_n})_{n \geq 0}$ with microlocal limit $\mu$, and that for any probability measure $\mu$ on $\Sigma$.

We now assume that $0 < E < E_c$. Fix $z \in \{p=E\}$ and let $(z_k)_{k \geq 0}$ be a sequence converging to $z$ such that $a(z_k) = k^{-1}m_k$ with $m_k \in \Z_{\geq 0}$. By Proposition~\ref{lem:gaussian-beam}, $\bmf_{k,m_k,z_k}/\|\bmf_{k,m_k,z_k}\|_{L^2}$ converges microlocally to $\delta_z$, the Dirac masses carried by the periodic orbit of $z$. As above, one can also realize any convex linear combination of such Dirac masses. As these are dense in flow-invariant probability measures on $\{p=E\}$, this proves the claim.

\smallskip

\eqref{theorem:mainii}
This follows immediately from Proposition \ref{proposition:scl-measure}, item \eqref{proposition:scl-measureii}, and the unique ergodicity of the horocyclic flow (see Theorem \ref{lemma:horocycle}, item \eqref{lemma:horocycleii}).

\smallskip

\eqref{theorem:mainiii}
The proof follows \emph{verbatim} the standard proof of Quantum Ergodicity. In~the framework of twisted quantization, this question was recently addressed in \cite[\S 5.3]{Cekic-Lefeuvre-24} (in a more general framework).
\end{proof}

\begin{thebibliography}{Cha24b}
\bibitem[Arn61]{Arnold-61} V. I. Arnol’d – “Some remarks on flows of line elements and frames”, Soviet Math. Dokl. 2 (1961), p. 562–564.
\bibitem[CL]{Cekic-Lefeuvre-24} M. Cekić \& T. Lefeuvre – “Semiclassical analysis on principal bundles”, arXiv:2405.14846.
\bibitem[Cha25]{Charles--23} L. Charles , “Resolvents of Bochner Laplacians in the semiclassical limit”, Comm. Partial Differential Equations 50 (2025), no. 7, p. 899–930.
\bibitem[Cha00]{Charles-00} L. Charles – “Aspects semi-classiques de la quantification géométrique”, PhD Thesis, Université Paris Dauphine-PSL, 2000.
\bibitem[Com86]{Comtet-86} A. Comtet – “On the Landau levels on the hyperbolic plane”, Ann. Physics 173 (1986), p. 185–209.
\bibitem[DJN22]{Dyatlov-Jin-Nonnenmacher-22} S. Dyatlov, L. Jin \& S. Nonnenmacher – “Control of eigenfunctions on surfaces of variable curvature”, J. Amer. Math. Soc. 35 (2022), no. 2, p. 361–465.
\bibitem[DZ19]{Dyatlov-Zworski-book} S. Dyatlov \& M. Zworski – Mathematical theory of scattering resonances, Grad. Stud. Math., vol. 200, American Mathematical Society, Providence, RI, 2019.
\bibitem[Gin96]{Ginzburg-96} V. L. Ginzburg – “On the existence and non-existence of closed trajectories for some Hamiltonian flows”, Math. Z. 223 (1996), no. 3, p. 397–409.
\bibitem[IL94]{Iengo-Li-94} R. Iengo \& D. Li – “Quantum mechanics and quantum Hall effect on Riemann surfaces”, Nuclear Phys. B 413 (1994), no. 3, p. 735–753.
\bibitem[Cha24a]{Landau1} L. Charles , “Landau levels on a compact manifold”, Ann. H. Lebesgue 7 (2024), p. 69–121.
\bibitem[Cha24b]{Landau2} L. Charles , “On the spectrum of nondegenerate magnetic Laplacians”, Anal. PDE 17 (2024), no. 6, p. 1907–1952.
\bibitem[Lef26]{Lefeuvre-book} T. Lefeuvre – Microlocal analysis in hyperbolic dynamics and geometry, Cours Spécialisés, vol. 32, Société Mathématique de France, Paris, 2026.
\bibitem[Non19]{Nonnenmacher-notes} S. Nonnenmacher – “Introduction to semiclassical analysis”, Lecture notes, available at https://www.imo.universite-paris-saclay.fr/~stephane.nonnenmacher/enseign/ Cours-Semiclassical-Analysis2019-lecture1-8.pdf, 2019.
\bibitem[TP06]{Tejero-06} C. Tejero Prieto – “Holomorphic spectral geometry of magnetic Schrödinger operators on Riemann surfaces”, Differential Geom. Appl. 24 (2006), no. 3, p. 288–310.
\bibitem[Wei77]{Weinstein-77} A. Weinstein – “Asymptotics of eigenvalue clusters for the Laplacian plus a potential”, Duke Math. J. 44 (1977), p. 883–892.
\bibitem[Zwo12]{Zworski-12} M. Zworski – Semiclassical analysis, Grad. Stud. Math., vol. 138, American Mathematical Society, Providence, RI, 2012.
\bibitem[Bur90]{Burger-90} M. Burger – “Horocycle flow on geometrically finite surfaces”, Duke Math. J. 61 (1990), no. 3, p. 779–803.
\bibitem[Fur73]{Furstenberg-73} H. Furstenberg – “The unique ergodicity of the horocycle flow”, in Recent advances in topological dynamics (New Haven, Conn., 1972), Lect. Notes in Math., vol. 318, Springer, Berlin-New York, 1973, p. 95–115.
\bibitem[FF03]{Flaminio-Forni-03} L. Flaminio \& G. Forni – “Invariant distributions and time averages for horocycle flows”, Duke Math. J. 119 (2003), no. 3, p. 465–526.
\bibitem[Pat99]{Paternain-99} G. P. Paternain – Geodesic flows, Progress in Math., vol. 180, Birkhäuser Boston, Inc., Boston, MA, 1999.
\end{thebibliography}
\end{document}
